% Part: first-order-logic
% Chapter: tableaux
% Section: derivations

\documentclass[../../../include/open-logic-section]{subfiles}

\begin{document}

\iftag{FOL}
      {\olfileid{fol}{tab}{der}}
      {\olfileid{pl}{tab}{der}}

\olsection{\usetoken{P}{tableau}}

\begin{explain}
قلنا ما الافتراض، وقدمنا قواعد الاستدلال. ومنهما تتولد
!!^{tableau}s استقرائيًا: فكل !!{tableau} إما أن يكون فرعًا واحدًا
يتألف من افتراض واحد أو أكثر، وإما أن ينتج من !!a{tableau} بتطبيق إحدى
قواعد الاستدلال على فرع.
\end{explain}

\begin{defn}[!!^{tableau}]
!!^a{tableau} للافتراضات~$\sFmla{S_1}{!A_1}$، \dots،
$\sFmla{S_n}{!A_n}$ (حيث يكون كل~$S_i$ إما~$\True$ أو~$\False$) شجرة
منتهية من ال!!{signed formula}s تحقق الشروط الآتية:
\begin{enumerate}
\item تكون ال!!{signed formula}s العليا وعددها~$n$ في الشجرة هي
  $\sFmla{S_i}{!A_i}$، واحدة تحت الأخرى.
\item تنتج كل !!{signed formula} في الشجرة ليست من الافتراضات عن تطبيق
  صحيح لقاعدة استدلال على !!a{signed formula} في الفرع فوقها.
\end{enumerate}
يكون فرع !!a{tableau} \emph{مغلقًا} إذا وفقط إذا احتوى على كل من
$\sFmla{\True}{!A}$ و$\sFmla{\False}{!A}$، ويكون \emph{مفتوحًا} فيما
عدا ذلك. ويكون !!^a{tableau} كل فرع فيه مغلقًا \emph{!!{tableau}
مغلقًا} (لمجموعة افتراضاته). وإذا لم يكن !!a{tableau} مغلقًا، أي إذا
احتوى على فرع مفتوح واحد على الأقل، كان \emph{مفتوحًا}.
\end{defn}

\begin{ex}
  كل مجموعة افتراضات بمفردها !!a{tableau}، لكنها لن تكون مغلقة بوجه
  عام. (ومن الواضح أنها لا تكون مغلقة إلا إذا احتوت الافتراضات أصلًا
  على زوج من ال!!{signed formula}s~$\sFmla{\True}{!A}$
  و$\sFmla{\False}{!A}$.)

  يمكننا أن نحصل من !!a{tableau} (مفتوح أو مغلق) على جدول جديد أكبر
  بتطبيق إحدى قواعد الاستدلال على !!a{signed formula}
  $\sFmla{S}{!A}$ فيه. وستلحق القاعدة !!{signed formula}s واحدة أو أكثر
  بنهاية كل فرع يحتوي على ورود~$\sFmla{S}{!A}$ الذي نطبق عليه القاعدة.

  اعتبر مثلًا الافتراض~$\sFmla{\True}{!A \land \lnot !A}$. وفيما يأتي
  ال!!{tableau} (المفتوح) الذي يتألف من ذلك الافتراض وحده:
  \begin{center}
    \begin{tableau}{}
      [\sFmla{\True}{\formula{A} \land \lnot \formula{A}}, just=\TAss]
    \end{tableau}{}
  \end{center}
  ونحصل منه على !!{tableau} جديد بتطبيق قاعدة~$\TRule{\True}{\land}$
  على الافتراض. وتتيح لنا تلك القاعدة إضافة سطرين جديدين إلى
  ال!!{tableau}، هما~$\sFmla{\True}{!A}$ و
  $\sFmla{\True}{\lnot !A}$:
  \begin{center}
    \begin{tableau}{}
      [\sFmla{\True}{\formula{A} \land \lnot \formula{A}}, just=\TAss
        [\sFmla{\True}{\formula{A}}, just={\TRule{\True}{\land}[1]},
          [\sFmla{\True}{\lnot \formula{A}}, just={\TRule{\True}{\land}[1]}
          ]
        ]
      ]
    \end{tableau}{}
  \end{center}
  عندما نكتب ال!!{tableau}s، نسجل القواعد التي طبقناها على اليمين
  (فمثلًا، تعني~$\TRule{\True}{\land} 1$ أن ال!!{signed formula} في ذلك
  السطر ناتجة عن تطبيق قاعدة~$\TRule{\True}{\land}$ على ال!!{signed formula}
  في السطر~$1$). ويحتوي هذا ال!!{tableau} الجديد الآن على
  !!{signed formula}s إضافية، لكن لا يمكننا تطبيق قاعدة إلا على واحدة
  منها ($\sFmla{\True}{\lnot !A}$)، وهي في هذه الحالة قاعدة
  $\TRule{\True}{\lnot}$. وينتج من ذلك ال!!{tableau} المغلق الآتي:
  \begin{center}
    \begin{tableau}{}
      [\sFmla{\True}{\formula{A} \land \lnot \formula{A}}, just=\TAss
        [\sFmla{\True}{\formula{A}}, just={\TRule{\True}{\land}[1]}
          [\sFmla{\True}{\lnot \formula{A}}, just={\TRule{\True}{\land}[1]}
            [\sFmla{\False}{\formula{A}}, just={\TRule{\True}{\lnot}[3]}, close]
          ]
        ]
      ]
    \end{tableau}{}
  \end{center}
\end{ex}

\end{document}
